\documentclass[a4paper]{article}

% Basic packages
% Español (México) con XeLaTeX: fontspec para fuentes Unicode, polyglossia
% para la división silábica y los nombres del documento en la variante mexicana.
\usepackage{fontspec}
\usepackage{polyglossia}
\setdefaultlanguage[variant=mexican]{spanish}
% Los nombres chinos van como «pinyin (original)»: xeCJK da a los caracteres
% Han su propia fuente; sin ella XeLaTeX los omite con un aviso.
% FandolSong no cubre algunos caracteres de nombres (U+73FA); WenQuanYi Zen Hei sí.
\usepackage[AutoFallBack=true]{xeCJK}
\setCJKmainfont{FandolSong-Regular.otf}
\setCJKfallbackfamilyfont{\CJKrmdefault}{WenQuanYi Zen Hei}
% polyglossia no traduce los nombres de listings.
\AtBeginDocument{\providecommand{\lstlistingname}{}\renewcommand{\lstlistingname}{Listado}%
  \providecommand{\lstlistlistingname}{}\renewcommand{\lstlistlistingname}{Índice de listados}}
\usepackage{amsmath, amssymb}
\usepackage{graphicx}
\usepackage[margin=2.5cm]{geometry}
\usepackage[most]{tcolorbox}
\usepackage{etoolbox}
\usepackage{listings}
\usepackage{hyperref}
\usepackage{booktabs}
\usepackage{subcaption}
\usepackage{float}

% Highlight boxes
\newtcolorbox{knowledgebox}[1]{
    enhanced, colback=blue!5!white, colframe=blue!75!black, colbacktitle=blue!75!black,
    coltitle=white, fonttitle=\bfseries, title={#1}, attach boxed title to top left={yshift=-2mm, xshift=2mm},
    boxrule=1pt, sharp corners
}

\newtcolorbox{importantbox}[1]{
    enhanced, colback=yellow!10!white, colframe=yellow!80!black, colbacktitle=yellow!80!black,
    coltitle=black, fonttitle=\bfseries, title={#1}, sharp corners
}

\newtcolorbox{warningbox}[1]{
    enhanced, colback=red!5!white, colframe=red!75!black, colbacktitle=red!75!black,
    coltitle=white, fonttitle=\bfseries, title={#1}, sharp corners
}

% Listings
\lstset{
    language=Python,
    basicstyle=\ttfamily\small,
    keywordstyle=\color{blue},
    stringstyle=\color{red!60!black},
    commentstyle=\color{green!60!black},
    breaklines=true,
    frame=single,
    numbers=left,
    numberstyle=\tiny\color{gray},
    captionpos=b,
    extendedchars=true
}

% Front-page metadata
\newcommand{\notetitle}{[LLM Agents SP25] Open Training Recipes: LLM Reasoning}
\newcommand{\noteauthors}{Elaboradas a partir de la clase de Hanna Hajishirzi}
\newcommand{\notedate}{2026-04-02}
\newcommand{\videochannel}{Berkeley RDI}
\newcommand{\videopublishdate}{2025-03-03}
\newcommand{\videoduration}{1:20:53}
\newcommand{\videourl}{https://www.youtube.com/watch?v=cMiu3A7YBks}
\newcommand{\videocoverpath}{cover.jpg}
\newcommand{\repourl}{https://github.com/hqhq1025/ai-course-notes}


% Footer with repo info
\usepackage{fancyhdr}
\pagestyle{fancy}
\fancyhf{}
\fancyfoot[C]{\thepage}
\fancyfoot[R]{\footnotesize\color{gray}\href{\repourl}{AI Course Notes}}
\renewcommand{\headrulewidth}{0pt}
\renewcommand{\footrulewidth}{0pt}

\begin{document}

\begin{titlepage}
\centering
{\Large Notas del curso\par}
\vspace{1.2cm}
{\huge\bfseries \notetitle\par}
\vspace{0.8cm}
{\large \noteauthors\par}
\vspace{0.3cm}
{\large \notedate\par}
\vspace{1.2cm}

\ifdefempty{\videocoverpath}{
{\small Al generar las notas, indique la ruta local de la portada del video.\par}
}{
\includegraphics[width=0.82\textwidth,height=0.45\textheight,keepaspectratio]{\videocoverpath}\par
}

\vfill
\begin{tcolorbox}[width=0.9\textwidth, colback=black!2!white, colframe=black!60, sharp corners]
\textbf{Autor o canal del video}: \videochannel\par
\textbf{Fecha de publicación}: \videopublishdate\par
\textbf{Duración del video}: \videoduration\par
\ifdefempty{\videourl}{}{
\textbf{Enlace del video}: \href{\videourl}{\nolinkurl{\videourl}}\par
}
\par
\textbf{Página del proyecto}: \href{\repourl}{\nolinkurl{\repourl}}\par
\end{tcolorbox}
\end{titlepage}

\tableofcontents
\newpage

%% --- Inicio del contenido --- %%

\section{Ubicación de la conferencia: por qué una ``receta de entrenamiento'' importa más que un truco puntual}

\subsection{Definición del problema: el razonamiento no es una capacidad única, sino un producto sistémico de la cadena de entrenamiento}
El valor central de esta conferencia no es presentar otro algoritmo nuevo, sino descomponer la ``capacidad de razonamiento'' hasta la realidad de la ingeniería: está determinada en conjunto por la distribución del preentrenamiento, los datos del posentrenamiento, el objetivo de optimización, el mecanismo de verificación y el presupuesto en tiempo de inferencia. Es decir, la capacidad de razonamiento no es el resultado directo de algún truco mágico, sino un problema de diseño de sistema de extremo a extremo.

\begin{importantbox}{Postura central del expositor}
\textit{``Open scientific research is why we are here.''} El expositor atribuye los avances de los LLM de los últimos años al ecosistema de investigación científica abierta, y enfatiza que cuando ese ecosistema se vuelve cerrado, disminuyen la reproducibilidad, la comparabilidad y la capacidad de corrección de la investigación.
\end{importantbox}

\begin{knowledgebox}{Por qué esta conferencia enfatiza las ``recipes''}
La discusión sobre la capacidad de los modelos suele reducirse a la escala de parámetros y al cómputo, pero una ``receta'' industrialmente reproducible en realidad contiene más variables implícitas:
\begin{itemize}
    \item la estrategia de muestreo y mezcla de datos;
    \item la división de las etapas de entrenamiento (preentrenamiento / posentrenamiento / tiempo de inferencia);
    \item el método de evaluación y la estrategia de descontaminación;
    \item la solución óptima de compromiso bajo restricciones de presupuesto.
\end{itemize}
\end{knowledgebox}

\begin{figure}[H]
\centering
\includegraphics[width=0.9\textwidth]{frame-01-open-science.jpg}
\caption{Video aproximadamente en 00:03:10: el expositor plantea explícitamente la ``investigación abierta'' como premisa de la metodología de la receta de entrenamiento}
\end{figure}

\subsection{De la ``publicación de modelos'' a la ``publicación científica''}
El expositor divide los modelos abiertos en tres niveles: solo pesos abiertos, código de entrenamiento abierto, y datos y proceso de evaluación abiertos. El que realmente sostiene la investigación científica es el tercer nivel, porque solo cuando se abren a la vez el procesamiento de datos, los scripts de entrenamiento y los criterios de evaluación puede la comunidad de investigación determinar de dónde proviene la mejora.

\begin{warningbox}{Abrir solo los pesos no equivale a ser reproducible}
Muchos trabajos solo publican el checkpoint, con lo cual los investigadores no pueden verificar si la mejora proviene de:
\begin{itemize}
    \item una mejora en la calidad de los datos;
    \item un ajuste fino de los hiperparámetros de entrenamiento;
    \item una modificación del optimizador o del programador;
    \item una contaminación de la evaluación o una diferencia de referencia.
\end{itemize}
Cuando estas variables no son visibles, el valor del llamado ``SOTA'' para la comunidad científica disminuye de manera notable.
\end{warningbox}

\subsection{Resumen de la sección}
Esta sección establece la línea base metodológica de toda la conferencia: la capacidad de razonamiento debe entenderse a partir de una ``receta de entrenamiento completa'', y la apertura y la reproducibilidad son condiciones necesarias para que esa receta pueda iterarse de manera acumulativa.

\section{Marco de tres etapas: Pre-training, Post-training, Inference-time}

\subsection{Perspectiva unificada: las tres palancas de la formación de capacidad}
El expositor divide el sistema de entrenamiento en tres palancas: el preentrenamiento determina la capacidad de la base y la cobertura de conocimiento, el posentrenamiento determina el estilo de comportamiento y las preferencias de tarea, y el escalamiento en tiempo de inferencia determina la cantidad de cómputo de búsqueda que se puede invertir por solicitud. Las tres son complementarias y presentan diferencias claras en el rendimiento marginal.

\begin{importantbox}{Las tres etapas no son una relación de sustitución, sino una relación de división del trabajo}
\begin{itemize}
    \item \textbf{Pre-training}: aporta una capacidad de representación de amplia cobertura y compresión de conocimiento;
    \item \textbf{Post-training}: orienta la capacidad hacia un comportamiento utilizable (conversación, seguimiento de instrucciones, alineación de preferencias);
    \item \textbf{Inference-time}: invierte cómputo adicional de manera dinámica en el momento del despliegue, para aumentar la tasa de éxito en tareas difíciles.
\end{itemize}
\end{importantbox}

\begin{figure}[H]
\centering
\includegraphics[width=0.9\textwidth]{frame-02-training-stack.jpg}
\caption{Video aproximadamente en 00:11:40: el expositor muestra la pila completa, desde el preentrenamiento hasta el posentrenamiento y el escalamiento en tiempo de prueba}
\end{figure}

\subsection{Descomposición de ingeniería: qué tipo de cuello de botella resuelve cada capa}
\begin{table}[H]
\centering
\renewcommand{\arraystretch}{1.2}
\begin{tabular}{p{2.6cm}p{3.2cm}p{4.0cm}p{3.2cm}}
\toprule
\textbf{Etapa} & \textbf{Objetivo principal} & \textbf{Técnicas típicas} & \textbf{Cuellos de botella comunes} \\
\midrule
Pre-training & Aprender representaciones, comprimir conocimiento, construir la capacidad de base & Proporción de datos, optimizador, programación de la tasa de aprendizaje, detalles de la arquitectura & Límite superior de la calidad de los datos, alto costo de cómputo \\
Post-training & Alinear el comportamiento, reforzar capacidades específicas & SFT, DPO/PPO, RLVR & Conflicto en la distribución de los datos, sesgo del reward \\
Inference-time & Aumentar la tasa de éxito y la robustez en problemas difíciles & Best-of-N, autoconsistencia, iteración de recuperación & Aumento de la latencia y el costo, amplificación de los errores de juicio \\
\bottomrule
\end{tabular}
\caption{Marco de entrenamiento de tres etapas y su división de cuellos de botella}
\end{table}

\begin{knowledgebox}{Por qué el curso enfatiza repetidamente la ``combinación''}
Una sola etapa difícilmente resuelve por sí sola el problema del razonamiento. Por ejemplo, si solo se hace posentrenamiento, se queda limitado por la capacidad de la base; si solo se hace escalamiento en tiempo de inferencia, se queda limitado por la calidad de los candidatos. Por eso, la ruta más realista es hacer una ``combinación acumulable'' según el presupuesto, en lugar de buscar un algoritmo universal.
\end{knowledgebox}

\subsection{Resumen de la sección}
El marco de tres etapas ofrece la perspectiva mínima y completa para analizar los modelos de razonamiento: primero se examina la base, luego la alineación, y por último la asignación de cómputo en el momento del despliegue.


\section{Preentrenamiento: la experiencia reproducible de OLMo 2}

\subsection{Objetivo: producir un modelo base muy competitivo dentro de un costo controlable}
El expositor usa OLMo 2 para discutir una realidad clave: los equipos de investigación no siempre cuentan con poder de cómputo de nivel frontier, pero aun así pueden, mediante una mejor gobernanza de datos y detalles de entrenamiento, producir una línea base sólida en las escalas de 7B/13B.

\begin{importantbox}{El valor de OLMo 2 está en la ``mejora interpretable''}
En comparación con solo publicar las curvas de resultados, la serie OLMo enfatiza más la publicación abierta de la cadena de entrenamiento, lo cual permite a la comunidad responder de manera sistemática:
\begin{itemize}
    \item qué detalles de entrenamiento realmente contribuyeron al desempeño;
    \item qué modificaciones solo son efectivas en puntos de referencia específicos;
    \item bajo distintos presupuestos, qué conjunto de configuraciones es más robusto.
\end{itemize}
\end{importantbox}

\begin{figure}[H]
\centering
\includegraphics[width=0.9\textwidth]{frame-03-olmo-scaling.jpg}
\caption{Video aproximadamente en 00:18:30: curvas de entrenamiento de OLMo y comparación de escalas, con énfasis en la optimización de la eficiencia bajo presupuesto restringido}
\end{figure}

\subsection{Los factores que realmente influyen en el modelo base de razonamiento durante el preentrenamiento}
\begin{knowledgebox}{Cuatro tipos de variables que se pasan por alto con más facilidad que la ``escala de parámetros''}
\begin{enumerate}
    \item la deduplicación de datos, el control de contaminación y la estratificación por calidad;
    \item la proporción de muestreo por dominio (código, matemáticas, texto de conocimiento, diálogo);
    \item la tasa de aprendizaje y la programación de warmup/cooldown;
    \item el tokenizer y la estrategia de ventana de contexto.
\end{enumerate}
\end{knowledgebox}

El expositor enfatiza que la capacidad de razonamiento ya queda parcialmente determinada en la etapa de preentrenamiento, y que el posentrenamiento solo puede reordenar y reforzar dentro de los límites de la capacidad ya existente. Esto explica por qué muchas técnicas de posentrenamiento producen resultados muy distintos según el modelo base.

\begin{warningbox}{Las métricas del preentrenamiento no corresponden una a una con la capacidad de razonamiento en tareas posteriores}
Un loss de modelado de lenguaje más bajo no necesariamente produce un razonamiento en cadena más sólido. Si el corpus de entrenamiento tiene señales insuficientes de estructura de razonamiento, el modelo puede ser hábil para continuar textos con fluidez, pero no para la inferencia con restricciones de varios pasos.
\end{warningbox}

\subsection{Resumen de la sección}
La lección de OLMo 2 es que, en un entorno de investigación, una receta de preentrenamiento abierta e interpretable constituye en sí misma un resultado, y determina el límite superior y la reproducibilidad de todas las mejoras de reasoning posteriores.


\section{Post-training I: SFT y ingeniería de datos}

\subsection{El papel del SFT: convertir capacidad latente en conducta invocable}
El SFT no es ``clases de regularización'', sino el establecimiento de un protocolo de interacción. Le indica al modelo base los límites de la tarea, el estilo de salida y la estructura de la respuesta, lo cual influye directamente en la estabilidad de la optimización de preferencias y del RL posteriores.

\begin{importantbox}{Juicio práctico del expositor}
Si los datos de SFT tienen una estructura desordenada, una cobertura de instrucciones estrecha y una calidad despareja, el DPO/PPO y el RLVR posteriores tenderán con más facilidad a presentar reward hacking, colapso de modos o un desarrollo desigual de las capacidades.
\end{importantbox}

\begin{figure}[H]
\centering
\includegraphics[width=0.9\textwidth]{frame-04-sft-data.jpg}
\caption{Video aproximadamente 00:29:30: fuentes de datos y estrategia de depuración en la etapa de SFT, con énfasis en la sinergia entre datos humanos y datos sintéticos}
\end{figure}

\subsection{Estrategia mixta de datos humanos y datos sintéticos}
\begin{knowledgebox}{El objetivo de la mezcla no es ``quién sustituye a quién'', sino ``quién complementa a quién''}
\begin{itemize}
    \item Datos anotados por humanos: de alta calidad y estilo genuino, pero de escala reducida y costo elevado;
    \item Datos sintéticos: de gran escala y amplia cobertura, pero con mayor riesgo de fluctuaciones de calidad y sesgo;
    \item La estrategia óptima suele consistir en anclar con muestras humanas de alta calidad y ampliar la cobertura con muestras sintéticas.
\end{itemize}
\end{knowledgebox}

El expositor mencionó en varias ocasiones la descontaminación de datos (decontamination) y la gobernanza de licencias, lo cual resulta especialmente crucial para la comunidad abierta, porque cuando existe contaminación de los benchmarks, las mejoras del entrenamiento pueden malinterpretarse como avances algorítmicos.

\begin{warningbox}{El error más común del SFT}
Fijarse solo en la cantidad total y no en la distribución. Acumular más datos de instrucción, si esto rompe el equilibrio entre tareas, puede hacer que la capacidad de razonamiento aumente mientras que la exactitud factual, el seguimiento del formato o la capacidad multilingüe disminuyan.
\end{warningbox}

\subsection{Resumen de la sección}
El núcleo de la etapa de SFT es la ingeniería de datos y no un algoritmo aislado. Un proceso de gobernanza de datos reproducible es la condición previa para que la optimización de preferencias y el aprendizaje por refuerzo posteriores funcionen de manera estable.


\section{Post-training II: optimización de preferencias (DPO/PPO)}

\subsection{De la señal supervisada a la señal de preferencias}
Después del SFT, el modelo ya cuenta con la capacidad básica para ejecutar tareas, pero el estilo de las respuestas, la preferencia de riesgo, la concisión y la utilidad todavía necesitan moldearse más mediante el aprendizaje de preferencias. El expositor presentó numerosas observaciones de ingeniería en torno a las diferencias prácticas entre DPO y PPO.

\begin{importantbox}{Por qué la optimización de preferencias es efectiva para reasoning}
Las tareas de reasoning suelen tener múltiples soluciones viables; la optimización de preferencias puede guiar al modelo hacia una distribución de salidas más interpretable, con pasos más completos y errores más corregibles, en lugar de limitarse a igualar una única etiqueta supervisada.
\end{importantbox}

\begin{figure}[H]
\centering
\includegraphics[width=0.9\textwidth]{frame-05-dpo-ppo.jpg}
\caption{Video aproximadamente en 00:48:30: comparación de DPO/PPO y discusión sobre la estabilidad del entrenamiento}
\end{figure}

\subsection{Las disyuntivas de ingeniería entre DPO y PPO}
\begin{table}[H]
\centering
\renewcommand{\arraystretch}{1.2}
\begin{tabular}{p{2.4cm}p{4.0cm}p{3.8cm}p{3.0cm}}
\toprule
\textbf{Método} & \textbf{Ventajas} & \textbf{Riesgos} & \textbf{Etapa de aplicación} \\
\midrule
DPO & La implementación es relativamente simple, el costo es menor y la cadena de entrenamiento es corta & Es sensible a la calidad de las muestras de preferencias y requiere manejar el sesgo de longitud & Iteración rápida a gran escala \\
PPO/RL en línea & Puede aprovechar la retroalimentación en línea para optimizar de manera continua, con mayor capacidad de exploración & Es inestable, el ajuste de hiperparámetros es complejo y el costo es alto & Etapas de impulso final para capacidades clave \\
\bottomrule
\end{tabular}
\caption{Diferencias prácticas entre DPO y PPO}
\end{table}

\begin{knowledgebox}{Un tema que se repite a lo largo del curso}
Cambiar de algoritmo casi nunca es el factor decisivo; \textbf{la calidad de la señal de recompensa} y \textbf{la cobertura de la distribución de entrenamiento} suelen ser más importantes. Muchas ``ganancias algorítmicas'' en realidad provienen de correcciones en los datos y en la función objetivo.
\end{knowledgebox}

\begin{warningbox}{Dos tipos de degradación implícita en el aprendizaje de preferencias}
\begin{itemize}
    \item \textbf{Complacencia excesiva}: el modelo aprende a «parecer una buena respuesta», pero carece de profundidad de razonamiento real;
    \item \textbf{Atajo de longitud}: el modelo alarga o acorta la salida para activar la preferencia de reward, en lugar de mejorar la corrección.
\end{itemize}
\end{warningbox}

\subsection{Resumen de la sección}
La optimización de preferencias es la capa de control del comportamiento dentro del post-training. DPO y PPO no son una disyuntiva binaria, sino que deben desplegarse por etapas combinando la calidad de los datos, el presupuesto y los objetivos de estabilidad.

\section{Post-training III: RLVR y recompensas verificables}

\subsection{Por qué RLVR se convirtió en el punto de inflexión clave para la mejora del razonamiento}
El expositor define RLVR (Reinforcement Learning with Verifiable Rewards) como el punto de inflexión clave dentro del post-training: cuando una tarea puede verificarse de manera automática como correcta o incorrecta, el modelo puede mejorar de forma continua a partir de una gran cantidad de muestras autogeneradas, sin depender ya por completo de anotaciones costosas de preferencia humana.

\begin{importantbox}{El ciclo mínimo de RLVR}
\begin{enumerate}
    \item Generar respuestas candidatas;
    \item Dar una recompensa verificable mediante un verificador basado en reglas o en un programa;
    \item Actualizar la política con RL para aumentar la probabilidad de las muestras con alta recompensa.
\end{enumerate}
\end{importantbox}

\begin{figure}[H]
\centering
\includegraphics[width=0.9\textwidth]{frame-06-rlvr.jpg}
\caption{Video aproximadamente en 01:03:20: transición del reward model neuronal hacia recompensas verificables, lo que reduce la dependencia de anotaciones}
\end{figure}

\subsection{Comparación entre modelos de recompensa neuronales y recompensas basadas en reglas}
\begin{knowledgebox}{Por qué el expositor enfatiza «verificable»}
Los modelos de recompensa neuronales son flexibles, pero presentan sesgos y una interpretabilidad limitada; las recompensas basadas en reglas son más limitadas, pero más confiables. Para tareas como matemáticas, código y razonamiento simbólico, las recompensas basadas en reglas pueden reducir de manera significativa el riesgo de reward hacking.
\end{knowledgebox}

\begin{table}[H]
\centering
\renewcommand{\arraystretch}{1.2}
\begin{tabular}{p{2.8cm}p{3.6cm}p{3.8cm}p{2.8cm}}
\toprule
\textbf{Tipo de recompensa} & \textbf{Ventajas} & \textbf{Desventajas} & \textbf{Tarea típica} \\
\midrule
Modelo de recompensa neuronal & Alta capacidad expresiva, puede cubrir objetivos subjetivos complejos & Propenso a sesgos, interpretabilidad débil, costo alto & Preferencias de diálogo, consistencia de estilo \\
Recompensa basada en reglas/verificable & Estable, auditable, permite entrenamiento a escala automatizado & Cobertura limitada, costo de diseño alto & Matemáticas, código, restricciones formalizadas \\
\bottomrule
\end{tabular}
\caption{Comparación entre los dos tipos de mecanismos de recompensa}
\end{table}

\begin{warningbox}{Condiciones límite de RLVR}
RLVR no mejora de manera automática todas las tareas. Si una tarea carece de un verificador confiable, o si el verificador está desalineado con el objetivo real, el entrenamiento empujará al modelo hacia la dirección errónea de «alta puntuación, baja calidad».
\end{warningbox}

\subsection{De la experiencia de DeepSeek-R1 a prácticas reproducibles en la comunidad abierta}
El expositor cita la inspiración de la ruta al estilo R1: cuando el verificador es económico y las muestras pueden generarse a escala, el RL puede amplificar rápidamente la capacidad de razonamiento. Para la comunidad abierta, lo importante es liberar en conjunto el verificador, los scripts de entrenamiento y el proceso de evaluación como código abierto; de lo contrario, es difícil verificar si las conclusiones son transferibles.

\subsection{Resumen de la sección}
El verdadero valor de RLVR consiste en transformar la «mejora del razonamiento», de una alta dependencia manual, en una «iteración programable». La condición previa es que la calidad del verificador sea confiable y esté alineada con el objetivo real.

\section{Inference-time：valor de ingeniería del Test-time Scaling}

\subsection{Idea central: destinar cómputo adicional a las muestras que más lo necesitan}
El presentador considera la ampliación en tiempo de prueba como una tercera palanca: no se modifican los parámetros, solo se agregan pasos de muestreo, verificación, recuperación o reflexión durante la inferencia, de modo que los problemas difíciles reciban un presupuesto de cómputo mayor. Esta estrategia suele ser más eficaz en el razonamiento de cadena larga.

\begin{importantbox}{Tres estrategias comunes de Test-time Scaling}
\begin{itemize}
    \item \textbf{Best-of-N / self-consistency}: se muestrea de forma diversa y luego se elige mediante reglas o mediante el propio modelo;
    \item \textbf{Iteración con recuperación aumentada}: se introduce evidencia externa a mitad del razonamiento para corregirlo de forma dinámica;
    \item \textbf{Búsqueda con restricciones de proceso}: se filtran las trayectorias intermedias mediante puntajes o restricciones de proceso.
\end{itemize}
\end{importantbox}

\begin{figure}[H]
\centering
\includegraphics[width=0.9\textwidth]{frame-07-testtime-scaling.jpg}
\caption{Video aproximadamente en 01:12:40: el beneficio y el costo de la ampliación en tiempo de prueba en tareas de razonamiento, en equilibrio}
\end{figure}

\subsection{Compensaciones en el despliegue: exactitud, latencia y costo}
\begin{table}[H]
\centering
\renewcommand{\arraystretch}{1.2}
\begin{tabular}{p{3.0cm}p{3.4cm}p{3.4cm}p{2.8cm}}
\toprule
\textbf{Estrategia} & \textbf{Cambio en la exactitud} & \textbf{Cambio en latencia/costo} & \textbf{Escenario adecuado} \\
\midrule
Decodificación única & Línea base & El más bajo & Escenarios en línea de alta concurrencia \\
Best-of-N & Mejora generalmente notable & Aumenta linealmente con N & Solicitudes de alto valor y baja frecuencia \\
Self-consistency + verificación & Gran beneficio en tareas verificables & Media a alta & Razonamiento matemático, de código o estructurado \\
Iteración con recuperación (Self-RAG) & Mejora la exactitud factual y la trazabilidad & Depende del costo de recuperación y reordenamiento & Preguntas y respuestas intensivas en conocimiento \\
\bottomrule
\end{tabular}
\caption{Compensaciones de despliegue de la ampliación en tiempo de prueba}
\end{table}

\begin{warningbox}{La ampliación en tiempo de prueba no es gratuita}
Si la calidad de los candidatos del modelo base es mala, muestrear más solo amplifica los errores; si el verificador no es confiable, el proceso de selección tomará el ruido como si fuera señal. Por ello, el límite superior del test-time scaling está doblemente restringido: por la calidad del modelo base y por el mecanismo de verificación.
\end{warningbox}

\begin{figure}[H]
\centering
\includegraphics[width=0.9\textwidth]{frame-08-selfrag.jpg}
\caption{Video aproximadamente en 01:17:20: los métodos al estilo Self-RAG plasman el ciclo cerrado de «razonamiento-recuperación-verificación»}
\end{figure}

\subsection{Resumen de la sección}
La ampliación en tiempo de prueba es una herramienta de asignación fina de presupuesto en el lado del despliegue, adecuada para solicitudes de alta dificultad y alto valor, pero debe usarse junto con un mecanismo de verificación confiable.


\section{Receta de extremo a extremo: la lógica de composición de OLMo a Tulu}

\subsection{Principio de composición: primero estabilizar la base, luego hacer la alineación y al final optimizar el presupuesto}
La ruta práctica real que propone esta clase no es «elegir primero el algoritmo», sino «elegir primero el orden de composición». Un orden ejecutable es:
\begin{enumerate}
    \item En la etapa de preentrenamiento, primero se establece una base estable;
    \item SFT ajusta el protocolo de la tarea y el comportamiento de salida a un rango controlable;
    \item DPO/PPO realiza el ajuste fino de preferencias;
    \item RLVR impulsa el razonamiento en subdominios verificables;
    \item el escalado en tiempo de prueba atiende las solicitudes de cola larga de alto valor.
\end{enumerate}

\begin{knowledgebox}{Por qué esta ruta es favorable para la comunidad abierta}
Cada capa puede experimentarse, evaluarse y liberarse como código abierto de forma independiente, lo que reduce la barrera de entrada para la colaboración en investigación y facilita localizar el problema cuando falla la reproducción de un experimento.
\end{knowledgebox}

\subsection{Sistema de evaluación: no basta con mirar un solo benchmark}
\begin{importantbox}{La perspectiva de evaluación de la clase}
La mejora de Reasoning debe verificarse conjuntamente en varias dimensiones:
\begin{itemize}
    \item tareas verificables como matemáticas y código;
    \item preguntas y respuestas de conocimiento y exactitud factual;
    \item estabilidad en contextos largos;
    \item restricciones de multilingüismo y seguridad;
    \item métricas de costo y latencia.
\end{itemize}
\end{importantbox}

\begin{warningbox}{El éxito en las métricas no equivale al éxito del sistema}
Si la evaluación no cubre las restricciones de despliegue (latencia, presupuesto, robustez), el modelo puede resultar inutilizable en un entorno de producto real aunque mejore en el benchmark.
\end{warningbox}

\subsection{Resumen de la sección}
El núcleo de la receta de extremo a extremo es el orden y la composición, no la optimización de un punto aislado. Los sistemas exitosos suelen ser el resultado de múltiples etapas «subóptimas pero complementarias».


\section{Riesgos y agenda de investigación: qué hacer a continuación en el entrenamiento abierto}

\subsection{Las tres brechas de investigación más críticas en la actualidad}
\begin{table}[H]
\centering
\renewcommand{\arraystretch}{1.2}
\begin{tabular}{p{3.0cm}p{4.0cm}p{4.2cm}}
\toprule
\textbf{Brecha} & \textbf{Problema actual} & \textbf{Dirección siguiente} \\
\midrule
Ecosistema de verificadores & La cobertura de tareas verificables sigue siendo limitada y la transferencia entre dominios es débil & Construir bibliotecas de verifiers componibles y estándares de estratificación de tareas \\
Gobernanza de datos & La calidad de los datos sintéticos es inestable y la contaminación es difícil de auditar de manera unificada & Unificar protocolos de descontaminación y cadenas de herramientas de auditoría de datos de código abierto \\
Control de costos & Los beneficios y la predicción de costos del escalamiento en tiempo de prueba son inestables & Establecer una estrategia de programación conjunta de ``calidad-latencia-costo'' \\
\bottomrule
\end{tabular}
\caption{Tres agendas de investigación aplicables derivadas del contenido de la clase}
\end{table}

\begin{knowledgebox}{Implicaciones directas para los equipos de ingeniería}
Documentar y versionar la ``receta de entrenamiento'', gestionándola con el mismo rigor que el código. Cada cambio de capacidad debe poder rastrearse hasta una modificación concreta en los datos, la función objetivo, el optimizador o la estrategia de inferencia.
\end{knowledgebox}

\subsection{Dos recomendaciones para investigadores y equipos de producto}
\begin{importantbox}{Recomendación 1: priorizar la construcción de una base experimental reproducible}
Aunque a corto plazo no se persiga la puntuación más alta, primero hay que dejar funcionando el entrenamiento, la evaluación y las pruebas de regresión. Sin una base reproducible, cualquier mejora puede resultar insostenible.
\end{importantbox}

\begin{importantbox}{Recomendación 2: tratar el presupuesto de inferencia como un parámetro de producto}
El escalamiento en tiempo de inferencia no es un accesorio algorítmico, sino una estrategia de producto. El presupuesto debe asignarse por niveles según el valor de la solicitud, en lugar de usar la misma estrategia de decodificación para todas las solicitudes.
\end{importantbox}

\begin{warningbox}{La apertura no significa estándares bajos}
La ruta del entrenamiento abierto también requiere un cumplimiento normativo de datos riguroso, descontaminación de la evaluación y revisión de seguridad. De lo contrario, la apertura terminará replicando ruido en lugar de replicar progreso.
\end{warningbox}

\subsection{Resumen de la sección}
El factor decisivo de la siguiente etapa del entrenamiento abierto no es un único algoritmo nuevo, sino esta infraestructura triple: el ecosistema de verificadores, la gobernanza de datos y la programación del presupuesto de inferencia.

\section{Manual de implementación: convertir la receta de entrenamiento en un proceso ejecutable para el equipo}

\subsection{Ejemplo de ritmo de iteración de cuatro semanas}
Para que el método de esta clase realmente se lleve a la práctica, lo más útil es convertir la receta de entrenamiento en un ritmo fijo, en lugar de experimentos improvisados. A continuación se presenta una plantilla de ritmo de cuatro semanas, utilizable en equipos de investigación de tamaño mediano:
\begin{enumerate}
    \item Semana 1: congelar la tarea objetivo y el protocolo de evaluación, y completar la descontaminación de datos y la ejecución de la línea base;
    \item Semana 2: experimentos de mezcla de datos de SFT, para confirmar el intervalo de estabilidad del comportamiento;
    \item Semana 3: barridos a pequeña escala de DPO/PPO, para verificar la ganancia de preferencia y los efectos secundarios;
    \item Semana 4: experimento conjunto de RLVR y test-time scaling, con el que se obtiene la curva de costo de despliegue.
\end{enumerate}

\begin{importantbox}{Hacer que el «ritmo experimental» importe más que la «inspiración»}
La mejora de la capacidad de razonamiento casi nunca es un salto único, sino la acumulación de pequeños cambios continuos. Solo al fijar los objetivos semanales, las métricas de evaluación y las pruebas de regresión se puede evitar que el equipo repita el mismo error una y otra vez.
\end{importantbox}

\subsection{Plantilla de registro de experimentos: campos mínimos rastreables}
\begin{table}[H]
\centering
\renewcommand{\arraystretch}{1.2}
\begin{tabular}{p{3.0cm}p{3.7cm}p{4.1cm}}
\toprule
\textbf{Campo} & \textbf{Contenido que se registra} & \textbf{Consecuencia de omitirlo} \\
\midrule
Versión de datos & Origen de las muestras, reglas de deduplicación, proporción de mezcla, estado de licencia & No se puede determinar si la ganancia proviene del algoritmo o de una deriva de los datos \\
Configuración de entrenamiento & Optimizador, tasa de aprendizaje, lote, número de pasos de entrenamiento, seed & Los resultados no son reproducibles y es difícil hacer una regresión \\
Definición de la recompensa & Reglas de preferencia/versión del verificador/umbral & Es difícil diagnosticar el reward hacking \\
Protocolo de evaluación & Conjunto de referencias, reglas de descontaminación, intervalo de confianza & Las puntuaciones no son comparables y las decisiones se distorsionan \\
Presupuesto de despliegue & Latencia promedio, P95, costo de tokens, tasa de reintento por fallas & No se puede determinar si es viable ponerlo en producción \\
\bottomrule
\end{tabular}
\caption{Campos mínimos de registro para institucionalizar la receta de entrenamiento}
\end{table}

\begin{knowledgebox}{El punto que más se pasa por alto en la ingeniería}
Muchos equipos solo registran el «mejor resultado» y no registran los «experimentos fallidos». Pero en el entrenamiento de reasoning, las trayectorias fallidas suelen revelar mejor que las trayectorias exitosas problemas de contaminación de datos, sesgo de recompensa o de la estrategia de muestreo.
\end{knowledgebox}

\subsection{De la investigación a la producción: lista de verificación del lanzamiento gradual}
\begin{warningbox}{Las cuatro verificaciones que deben completarse antes de salir a producción}
\begin{itemize}
    \item \textbf{Lanzamiento gradual de corrección}: limitar primero el tráfico en tareas de alto valor, para verificar la tasa de aprobación real;
    \item \textbf{Lanzamiento gradual de costo}: comparar el beneficio/costo de la decodificación única frente al test-time scaling;
    \item \textbf{Lanzamiento gradual de estabilidad}: observar si las salidas de cadena larga presentan deriva de formato o razonamiento repetido;
    \item \textbf{Lanzamiento gradual de seguridad}: realizar pruebas de regresión de equipo rojo sobre prompts de alto riesgo y solicitudes que exceden los permisos autorizados.
\end{itemize}
\end{warningbox}

\subsection{Resumen de la sección}
Que la receta de entrenamiento genere valor a largo plazo depende de si se institucionaliza como un proceso del equipo. Un ritmo fijo, un registro completo y las verificaciones del lanzamiento gradual son lo que permite convertir los resultados de investigación en una capacidad de producto estable.

\section{Desglose de un caso: cómo diseñar un experimento de mejora de reasoning}

\subsection{Objetivos y restricciones del experimento}
Supongamos que queremos mejorar un modelo de la escala 7B en 5 a 10 puntos en tareas de razonamiento matemático y de código, sin dañar de manera notable la capacidad de conversación general. La restricción de presupuesto es ``entrenamiento controlable, inferencia lista para producción''. Este tipo de problema corresponde exactamente a la ruta combinada que propone esta clase.

\begin{importantbox}{Cómo definir los objetivos}
El objetivo debe incluir al mismo tiempo:
\begin{itemize}
    \item \textbf{Objetivo de capacidad}: mejora de la exactitud en tareas verificables;
    \item \textbf{Límite de efectos secundarios}: la caída en tareas generales no debe superar un umbral;
    \item \textbf{Objetivo de despliegue}: costo y latencia dentro de un rango aceptable.
\end{itemize}
Si falta cualquiera de estos elementos, la dirección de la optimización se desvía de la necesidad real.
\end{importantbox}

\subsection{Matriz de experimentos por etapas}
\begin{table}[H]
\centering
\renewcommand{\arraystretch}{1.2}
\begin{tabular}{p{2.8cm}p{3.4cm}p{3.8cm}p{3.0cm}}
\toprule
\textbf{Etapa} & \textbf{Variable} & \textbf{Métrica observada} & \textbf{Condición de detención} \\
\midrule
SFT & Proporción de datos humanos/sintéticos, distribución de longitud & Seguimiento de instrucciones, estabilidad de formato, exactitud básica & Revertir en cuanto el efecto secundario supere el umbral \\
DPO/PPO & Peso de la pérdida de preferencia, temperatura de muestreo & Utilidad, verbosidad, tasa de rechazo de respuesta & El beneficio marginal de la preferencia disminuye \\
RLVR & Reglas del verificador, número de rondas de muestreo, número de pasos de actualización & Tasa de aprobación verificable en matemáticas/código & Aparece una señal de atajo de recompensa \\
Inference-time & Valor de N, estrategia de reordenamiento, número de rondas de recuperación & Tasa de aprobación, latencia P95, costo por solicitud & El costo supera el presupuesto \\
\bottomrule
\end{tabular}
\caption{Matriz por etapas de un experimento de mejora de reasoning}
\end{table}

\subsection{Prioridad del análisis de errores}
\begin{knowledgebox}{Los tres tipos de muestras de error más valiosos}
\begin{enumerate}
    \item \textbf{Error de alta confianza}: el modelo produce un razonamiento largo y encadenado, expresado con confianza pero incorrecto;
    \item \textbf{Juicio erróneo del verificador}: la respuesta es correcta pero la regla la descarta por error, o al revés;
    \item \textbf{Muestras de presupuesto ineficiente}: requieren un muestreo masivo para lograr una mejora mínima en la exactitud.
\end{enumerate}
Corregir primero estas muestras suele mejorar al mismo tiempo la calidad y la eficiencia de costos.
\end{knowledgebox}

\begin{warningbox}{No ocultar fallas estructurales con el puntaje promedio}
Un aumento en el puntaje promedio puede ocultar el retroceso en tareas clave. Un sistema real necesita examinar más bien ``el desempeño por segmento de tareas'': por ejemplo, cómo cambian por separado las matemáticas de varios pasos, la reparación de código largo y las preguntas y respuestas que dependen de recuperación.
\end{warningbox}

\subsection{Cómo relacionar este caso con el hilo principal de la clase}
Este caso muestra que el marco de tres etapas y la idea de receta abierta de esta clase pueden traducirse directamente en decisiones de ingeniería:
\begin{itemize}
    \item El preentrenamiento determina el límite superior del modelo base;
    \item El posentrenamiento determina la dirección del comportamiento;
    \item El escalamiento en tiempo de inferencia determina la curva de beneficio en el lado del despliegue;
    \item La apertura y la reproducibilidad determinan si el equipo puede iterar de manera sostenida.
\end{itemize}

\subsection{Resumen de la sección}
Un experimento exitoso de mejora de reasoning no es la victoria de un solo parámetro puntual, sino el resultado de que las restricciones del objetivo, la matriz de etapas, el análisis de errores y la gestión del presupuesto trabajen de manera coordinada.

\section{Síntesis y ampliación}

\subsection{Tabla de síntesis de toda la clase}
\begin{table}[H]
\centering
\renewcommand{\arraystretch}{1.2}
\begin{tabular}{p{2.7cm}p{3.9cm}p{4.2cm}}
\toprule
\textbf{Etapa} & \textbf{Pregunta central} & \textbf{Respuesta ejecutable que ofrece esta clase} \\
\midrule
Preentrenamiento & Cómo fortalecer la capacidad del modelo base dentro de un presupuesto & Usar la receta reproducible de OLMo para optimizar los datos y los detalles de entrenamiento \\
Post-training & Cómo convertir la capacidad en un comportamiento utilizable y una ventaja de razonamiento & SFT establece el protocolo, DPO/PPO ajusta las preferencias, RLVR impulsa el razonamiento verificable \\
Inference-time & Cómo mejorar aún más la eficiencia en el extremo de despliegue & Usar test-time scaling y el ciclo cerrado de verificación mediante recuperación según el valor de cada solicitud \\
Ecosistema abierto & Cómo hacer que las mejoras se puedan acumular & Abrir los datos, el código y la cadena de evaluación, y reforzar la reproducibilidad y la capacidad de auditoría \\
\bottomrule
\end{tabular}
\caption{Conclusión estructurada de Open Training Recipes for Reasoning}
\end{table}

\subsection{Conclusión central}
\begin{importantbox}{Resumen en una frase}
La capacidad de Reasoning es resultado de la ingeniería de sistemas, no de una técnica puntual; la ruta más eficaz es la receta combinada de ``preentrenamiento reproducible + posentrenamiento por etapas + escalamiento presupuestado en tiempo de prueba''.
\end{importantbox}

\begin{knowledgebox}{Los tres puntos de esta clase más valiosos para la práctica}
\begin{enumerate}
    \item Primero construir una cadena experimental reproducible, y después buscar puntuaciones más altas;
    \item Usar recompensas verificables para ampliar el alcance de automatización del entrenamiento por RL;
    \item Tratar el test-time scaling como una capacidad de gestión del presupuesto de despliegue, y no como un parche temporal.
\end{enumerate}
\end{knowledgebox}

\subsection{Lecturas adicionales}
\begin{itemize}
    \item Informes técnicos de OLMo / OLMo 2 y el código de entrenamiento (AI2)
    \item El artículo sobre la receta de posentrenamiento de Tulu 3 y su implementación de código abierto
    \item Los artículos de la serie DPO / PPO sobre optimización de preferencias y los métodos de corrección del sesgo de longitud
    \item Material público relacionado con RLVR y la ruta de DeepSeek-R1
    \item Investigación relacionada con Self-RAG y test-time scaling
\end{itemize}

\subsection{Preguntas de ampliación}
\begin{warningbox}{Preguntas abiertas que vale la pena seguir de cerca}
\begin{itemize}
    \item ¿Cómo construir verificadores confiables y de bajo costo para tareas que no son de matemáticas ni de código?
    \item ¿Cómo establecer un protocolo de evaluación consistente entre modelos y entre versiones de datos?
    \item ¿Cómo decidir de forma dinámica, con un presupuesto fijo, si vale la pena ``volver a muestrear''?
\end{itemize}
\end{warningbox}

%% --- Fin del contenido --- %%

\end{document}
